% =========================================================
% Stage 0 S0-B
% =========================================================
\chapter{S0-B --- Basic Containers \& STL}

% =========================================================
% Part A
% =========================================================
\section{Part A --- 위치와 선수지식}

\subsection{현재 위치}
S0-A에서 표준 입력과 출력, 기본 자료형, 조건문, 반복문, 함수, scalar 상태 전이를 다뤘다면 S0-B에서는 여러 값을 저장하고 순회하고 정렬하기 위한 C++ 표준 container와 STL의 기초를 다룬다.

이번 Learning Unit의 핵심 범위는 다음과 같다.
\begin{itemize}
  \item \code{string}
  \item \code{array}
  \item \code{vector}
  \item \code{pair}
  \item index loop, range-based \code{for}, iterator
  \item \code{begin()}, \code{end()}
  \item \code{sort()}, \code{reverse()}
  \item custom comparator
  \item pass-by-value, reference, \code{const} reference
\end{itemize}

\subsection{왜 지금 배우는가}
Competitive Programming에서는 입력이 하나의 값이 아니라 배열, 문자열, 레코드 목록으로 주어지는 경우가 많다. 이때 중요한 것은 단순히 container 문법을 외우는 것이 아니라 다음을 스스로 판단하는 것이다.
\begin{itemize}
  \item 고정 크기와 가변 크기 중 어느 형태가 필요한가?
  \item 원소의 index가 필요한가, 값만 순회하면 되는가?
  \item 정렬 기본 순서로 충분한가, custom comparator가 필요한가?
  \item 함수에 container를 넘길 때 복사해야 하는가, 원본을 수정해야 하는가, 읽기만 하면 되는가?
\end{itemize}

이 판단은 이후 two pointers, sliding window, binary search, graph, DP 등 거의 모든 알고리즘 구현에서 반복해서 사용된다.

\subsection{Core Decision Boundary}
S0-B에서는 다음 선택 기준을 핵심 경계로 본다.
\begin{center}
\begin{tabularx}{\textwidth}{p{4.0cm}X}
\toprule
선택 상황 & 판단 기준 \\
\midrule
\code{array} vs \code{vector} & 컴파일 시점에 고정된 크기인가, 실행 중 정해지거나 변할 수 있는 크기인가? \\
index loop vs range-for & index가 필요한가, 값 자체만 순회하면 되는가? \\
기본 \code{sort} vs comparator & 기본 오름차순/사전순으로 충분한가, 여러 우선순위 규칙이 필요한가? \\
값 전달 vs reference & 복사가 필요한가, 원본을 수정해야 하는가, 읽기만 하면서 복사는 피하고 싶은가? \\
\bottomrule
\end{tabularx}
\end{center}

\subsection{Immediate Coverage Floor}
이 Unit을 즉시 통과하기 위한 최소 증거는 다음과 같이 잡는다.
\begin{itemize}
  \item \code{vector}와 \code{string}을 이용한 독립 구현
  \item 기본 STL 알고리즘을 실제 코드에서 사용
  \item 최소 하나의 선택 이유 설명: \code{array}/\code{vector}, 기본 정렬/comparator, 값/reference 중 하나 이상
\end{itemize}

% =========================================================
% Part B
% =========================================================
\section{Part B --- 학습 목표}
이 Unit을 마치면 다음 행동을 독립적으로 수행할 수 있어야 한다.

\begin{enumerate}
  \item \code{string}, \code{array}, \code{vector}, \code{pair}의 역할을 구분한다.
  \item \code{vector}에 원소를 추가하고 index 또는 iterator로 접근한다.
  \item \code{pair}의 \code{first}, \code{second}를 이용해 두 값을 하나의 record로 묶는다.
  \item index loop, range-based \code{for}, iterator loop 중 상황에 맞는 순회를 선택한다.
  \item \code{begin()}과 \code{end()}가 나타내는 iterator 범위를 설명한다.
  \item \code{sort()}와 \code{reverse()}를 사용할 수 있다.
  \item 기본 정렬 규칙과 다른 우선순위가 필요한 경우 comparator를 작성한다.
  \item comparator가 strict ordering을 표현해야 한다는 점을 이해한다.
  \item \code{T}, \code{T\&}, \code{const T\&} 전달 방식의 차이를 설명한다.
  \item container나 문자열을 불필요하게 복사하지 않도록 함수 매개변수를 설계한다.
  \item \code{sort}의 기본 시간복잡도 $O(N\log N)$과 선형 순회의 $O(N)$을 구분한다.
  \item signed/unsigned 비교와 iterator invalidation 같은 기본 안전성 이슈를 인식한다.
\end{enumerate}

% =========================================================
% Part C
% =========================================================
\section{Part C --- 개념과 원리}

\subsection{C1. \code{string}: 문자 하나가 아니라 문자열 container}
\code{string}은 문자들의 연속된 sequence를 다루는 표준 타입이다.
\begin{lstlisting}
string s = "apple";
cout << s.size();       // 5
cout << s[0];           // 'a'
s.push_back('!');
\end{lstlisting}

\code{cin >> s;}는 공백 전까지의 한 token을 읽는다. 공백을 포함한 한 줄 전체를 읽어야 한다면 \code{getline}을 사용한다.

\begin{notebox}[문자열 비교]
\code{s1 < s2}는 기본적으로 사전식(lexicographic) 비교를 한다. 문자열 길이 자체를 먼저 비교하고 싶다면 \code{s.size()}를 별도로 비교해야 한다.
\end{notebox}

\subsection{C2. \code{array}: 크기가 타입의 일부인 고정 크기 container}
\begin{lstlisting}
array<int, 5> a = {1, 2, 3, 4, 5};
\end{lstlisting}

\code{array<int,5>}와 \code{array<int,10>}은 서로 다른 타입이다. 따라서 크기가 컴파일 시점에 고정되어 있고 바뀌지 않는 경우에 적합하다.

\subsection{C3. \code{vector}: 실행 중 크기를 정할 수 있는 연속 container}
\begin{lstlisting}
vector<int> v;
v.push_back(10);
v.push_back(20);
\end{lstlisting}

또는 처음부터 크기를 정할 수 있다.
\begin{lstlisting}
vector<int> v(100);
\end{lstlisting}

initializer list도 사용할 수 있다.
\begin{lstlisting}
vector<int> v = {3, 1, 4, 1, 5};
\end{lstlisting}

\code{push\_back()}은 평균적으로 amortized $O(1)$이다. 다만 내부 저장 공간이 부족해지면 더 큰 메모리 공간으로 재할당(reallocation)할 수 있으며, 이때 기존 원소를 가리키던 iterator나 reference가 무효화될 수 있다.

\subsection{C4. \code{pair}: 두 값을 하나의 record로 묶기}
\begin{lstlisting}
pair<string, int> p = {"apple", 5};
cout << p.first;   // apple
cout << p.second;  // 5
\end{lstlisting}

여러 record가 필요하다면 다음처럼 쓸 수 있다.
\begin{lstlisting}
vector<pair<string, int>> records;
records.push_back({"alice", 90});
records.push_back({"bob", 80});
\end{lstlisting}

C++17에서는 structured binding도 사용할 수 있다.
\begin{lstlisting}
for (const auto& [name, score] : records) {
    cout << name << ' ' << score << '\n';
}
\end{lstlisting}

\subsection{C5. 세 가지 대표 순회 방식}
index가 필요하다면 index loop가 자연스럽다.
\begin{lstlisting}
for (size_t i = 0; i < v.size(); ++i) {
    cout << i << ' ' << v[i] << '\n';
}
\end{lstlisting}

값만 필요하다면 range-based \code{for}가 간결하다.
\begin{lstlisting}
for (const auto& x : v) {
    cout << x << '\n';
}
\end{lstlisting}

iterator를 직접 다루어야 한다면 다음과 같이 쓴다.
\begin{lstlisting}
for (auto it = v.begin(); it != v.end(); ++it) {
    cout << *it << '\n';
}
\end{lstlisting}

\subsection{C6. \code{begin()}과 \code{end()}: 원소도 index도 아닌 iterator}
다음 vector를 생각하자.
\begin{lstlisting}
vector<int> v = {10, 20, 30, 40};
\end{lstlisting}

\code{v.begin()}은 첫 번째 원소를 가리키는 iterator이고, \code{v.end()}는 마지막 원소 바로 다음 위치를 나타내는 iterator다.

\begin{center}
\begin{tabular}{c|cccc|c}
index & 0 & 1 & 2 & 3 & 4에 해당하는 위치 \\
\hline
값 & 10 & 20 & 30 & 40 & 실제 원소 없음 \\
iterator & \code{begin()} & & & & \code{end()} \\
\end{tabular}
\end{center}

iterator가 가리키는 실제 원소를 얻으려면 dereference 연산자 \code{*}를 사용한다.
\begin{lstlisting}
auto it = v.begin();
cout << *it;  // 10
\end{lstlisting}

\code{v.end()}는 실제 원소를 가리키지 않으므로 다음 코드는 잘못된 사용이다.
\begin{lstlisting}
cout << *v.end();  // invalid
\end{lstlisting}

STL의 범위는 일반적으로 반열린 구간 $[\text{first},\text{last})$를 사용한다. 따라서
\begin{lstlisting}
sort(v.begin(), v.end());
\end{lstlisting}
는 첫 원소부터 \code{end()} 직전, 즉 vector 전체를 정렬한다.

\subsection{C7. \code{auto}: 초기값으로부터 타입 추론}
\code{auto}는 타입이 없는 변수를 만드는 키워드가 아니다. 컴파일러가 초기값을 보고 정적 타입을 추론하게 한다.
\begin{lstlisting}
auto x = 10;       // int
auto y = 3.14;     // double
auto c = 'a';      // char
auto it = v.begin();
\end{lstlisting}

마지막 줄에서 실제 iterator 타입을 길게 적는 대신 \code{auto}를 쓰는 것이 특히 편리하다.

다만 \code{auto}는 기본적으로 값 타입을 추론하므로 reference를 유지하고 싶으면 \code{\&}를 함께 써야 한다.
\begin{lstlisting}
for (auto x : v) {        // copy
    // ...
}

for (auto& x : v) {       // modifiable reference
    // ...
}

for (const auto& x : v) { // read-only reference
    // ...
}
\end{lstlisting}

\subsection{C8. reference: 기존 객체의 또 다른 이름}
다음을 보자.
\begin{lstlisting}
int x = 10;
int& r = x;
\end{lstlisting}

\code{r}은 \code{x}의 값을 복사한 새 \code{int}가 아니다. 이미 존재하는 \code{x} 객체에 대한 reference, 즉 alias다.
\begin{lstlisting}
r = 20;
cout << x;  // 20
\end{lstlisting}

반면 다음은 복사다.
\begin{lstlisting}
int x = 10;
int r = x;
r = 20;
cout << x;  // 10
\end{lstlisting}

\begin{keybox}[reference 선언 읽기]
\code{int\& r = x;}는 ``\code{r}을 \code{int} 객체 \code{x}에 대한 reference로 선언하고 \code{x}에 bind한다''라고 읽는 것이 정확하다.
\end{keybox}

reference는 선언할 때 참조 대상을 정해야 하며, 한 번 bind된 reference를 나중에 다른 객체의 reference로 다시 바꾸는 것은 아니다.

\subsection{C9. 함수 매개변수에서의 값, reference, const}
다음 세 선언을 비교하자.
\begin{lstlisting}
void f1(int& x) { ... }
void f2(int x) { ... }
void f3(const int x) { ... }
\end{lstlisting}

자연어로 읽으면 다음과 같다.
\begin{itemize}
  \item \code{int\& x}: 호출자가 넘긴 원본 \code{int}를 직접 참조하며, 함수 내부에서 원본을 수정할 수 있다.
  \item \code{int x}: 전달받은 값을 복사하여 함수의 지역 변수 \code{x}로 사용한다.
  \item \code{const int x}: 전달받은 값을 복사하되 함수 내부에서 그 복사본을 수정하지 못한다.
\end{itemize}

큰 객체를 읽기만 할 때는 다음 형태가 중요하다.
\begin{lstlisting}
void process(const vector<int>& v) {
    // v is not copied and cannot be modified through this reference
}
\end{lstlisting}

이는 ``vector 전체를 복사하지 않고 원본을 참조하되, 이 함수에서는 수정하지 않겠다''라는 뜻이다.

\subsection{C10. 일반 reference와 reference parameter의 관계}
다음 두 코드에서 \code{int\& a} 자체의 타입 의미는 같다.
\begin{lstlisting}
int& a = x;
\end{lstlisting}
\begin{lstlisting}
void func(int& a) {
    // ...
}
\end{lstlisting}

첫 번째는 선언하는 순간 \code{x}에 bind된다. 두 번째는 함수가 호출될 때 실제 argument에 bind되는 reference parameter다.
\begin{lstlisting}
int x = 10;
int y = 20;

func(x);  // this call: a refers to x
func(y);  // this call: a refers to y
\end{lstlisting}

따라서 reference parameter는 ``각 함수 호출에서 전달된 원본 객체에 붙는 alias''라고 이해할 수 있다.

\subsection{C11. \code{sort()}와 기본 정렬}
\code{sort(first,last)}는 $[\text{first},\text{last})$ 범위를 오름차순으로 정렬한다.
\begin{lstlisting}
vector<int> v = {4, 1, 3, 2};
sort(v.begin(), v.end());
// {1, 2, 3, 4}
\end{lstlisting}

일반적인 comparison sort의 시간복잡도는 $O(N\log N)$이다.

\code{pair}는 기본적으로 첫 번째 원소를 먼저 비교하고, 같으면 두 번째 원소를 비교하는 lexicographic ordering을 갖는다.

\subsection{C12. custom comparator}
기본 순서가 아닌 정렬이 필요하면 comparator를 전달한다.
\begin{lstlisting}
bool cmp(const pair<string, int>& a,
         const pair<string, int>& b)
{
    if (a.second != b.second) {
        return a.second > b.second;
    }
    return a.first < b.first;
}
\end{lstlisting}

\code{cmp(a,b)}의 의미는 ``\code{a}가 \code{b}보다 앞에 와야 하는가?''이다.

comparator는 strict ordering을 표현해야 한다. 같은 두 원소에 대해 \code{cmp(a,a)}가 참이 되는 \code{<=} 형태는 피해야 한다.

\subsection{C13. comparator 매개변수에서 \code{const \&}를 쓰는 이유}
네 형태를 비교하면 의미가 분명해진다.
\begin{lstlisting}
bool cmp(string s1, string s2);             // copies strings
bool cmp(string& s1, string& s2);           // no copy, may modify
bool cmp(const string s1, const string s2); // copies, local copies const
bool cmp(const string& s1, const string& s2); // no copy, read-only
\end{lstlisting}

문자열이나 record처럼 복사 비용이 있을 수 있는 객체를 비교만 할 때는 마지막 형태가 가장 자연스럽다.

\begin{keybox}[기억할 기준]
\code{\&}는 ``복사 대신 원본을 참조하는가''와 관련되고, \code{const}는 ``그 reference를 통해 수정할 수 있는가''와 관련된다.
\end{keybox}

\subsection{C14. \code{int}, \code{long long}, \code{size\_t}}
C++에는 여러 정수형이 있지만 CP에서 자주 만나는 대표적인 세 타입은 다음과 같다.
\begin{center}
\begin{tabular}{llll}
\toprule
타입 & signed 여부 & 전형적 크기 & 대표 용도 \\
\midrule
\code{int} & signed & 보통 32-bit & 일반 정수, 작은 index/개수 \\
\code{long long} & signed & 최소 64-bit & 큰 합, 큰 곱, 누적값 \\
\code{size\_t} & unsigned & 보통 64-bit 환경에서 64-bit & container 크기, 메모리 크기 \\
\bottomrule
\end{tabular}
\end{center}

\code{size\_t}는 ``양수라서 쓰는 타입''이라기보다 container의 크기나 메모리 크기를 표현하는 unsigned integer type이다. \code{vector::size()}와 \code{string::size()}가 반환하는 타입 계열도 여기에 해당한다.

\subsection{C15. signed/unsigned 비교 경고}
다음 코드를 보자.
\begin{lstlisting}
for (int j = 0; j < v.size(); ++j) {
    // ...
}
\end{lstlisting}

\code{j}는 signed \code{int}, \code{v.size()}는 unsigned \code{size\_t} 계열이므로 compiler가 signed/unsigned 비교 경고를 낼 수 있다.

정방향으로만 증가하고 $N$이 작은 문제에서는 실제 오류가 발생하지 않을 수도 있다. 그러나 일반적으로 signed 음수가 unsigned와 비교되면 큰 양수로 변환될 수 있다.

예를 들어 개념적으로 \code{-1}이 64-bit unsigned로 변환되면 매우 큰 양수가 될 수 있다. 따라서 다음과 같은 비교는 사람의 직관과 다른 결과를 만들 수 있다.
\begin{lstlisting}
int j = -1;
size_t n = 10;
// j < n may not behave like mathematical -1 < 10
\end{lstlisting}

index가 필요 없는 출력이라면 range-based \code{for}가 더 자연스러울 수 있다.
\begin{lstlisting}
for (const string& s : v) {
    cout << s << '\n';
}
\end{lstlisting}

\subsection{C16. unsigned 역순 순회의 함정}
다음 형태는 위험하다.
\begin{lstlisting}
for (size_t i = v.size() - 1; i >= 0; --i) {
    // ...
}
\end{lstlisting}

\code{size\_t}는 unsigned이므로 \code{i==0}에서 감소시키면 음수 \code{-1}이 아니라 매우 큰 값으로 underflow한다. 또한 unsigned 값에 대해 \code{i >= 0}은 항상 참이다.

역순 순회에서는 이러한 unsigned 특성을 의식해야 한다.

% =========================================================
% Part D
% =========================================================
\section{Part D --- Worked Example}

\subsection{D1. 학생 점수 record 정렬}
다음 목표를 생각하자.
\begin{quote}
학생의 이름과 점수를 입력받아 점수가 높은 학생부터 출력한다. 점수가 같으면 이름이 사전순으로 빠른 학생을 먼저 출력한다.
\end{quote}

두 값을 한 record로 묶기 위해 \code{pair<string,int>}를 사용하고, 여러 record를 저장하기 위해 \code{vector}를 사용한다.

\begin{lstlisting}
#include <algorithm>
#include <iostream>
#include <string>
#include <vector>
using namespace std;

bool cmp(const pair<string, int>& a,
         const pair<string, int>& b)
{
    if (a.second != b.second) {
        return a.second > b.second;
    }
    return a.first < b.first;
}

int main()
{
    int N;
    cin >> N;

    vector<pair<string, int>> students;

    for (int i = 0; i < N; ++i) {
        string name;
        int score;
        cin >> name >> score;
        students.push_back({name, score});
    }

    sort(students.begin(), students.end(), cmp);

    for (const auto& student : students) {
        cout << student.first << ' '
             << student.second << '\n';
    }
}
\end{lstlisting}

\subsection{D2. 이 예제에서 확인할 것}
\begin{itemize}
  \item \code{vector<pair<string,int>>}는 여러 record를 순서대로 저장한다.
  \item comparator는 record를 복사하지 않도록 \code{const \&}로 받는다.
  \item \code{sort(students.begin(), students.end(), cmp)}에서 \code{begin()}과 \code{end()}는 전체 정렬 범위를 나타낸다.
  \item 출력에서 index가 필요하지 않으므로 range-based \code{for}가 자연스럽다.
  \item 전체 시간복잡도는 정렬이 지배하므로 $O(N\log N)$이다. 문자열 비교 비용까지 일반화하면 문자열 길이도 고려할 수 있다.
\end{itemize}

% =========================================================
% Part E
% =========================================================
\section{Part E --- 중간평가 문제}

\subsection{E-1. Word Catalog}
\begin{problembox}{중간평가 문제 E-1 --- Word Catalog}
$N$개의 영어 소문자 문자열이 주어진다. 모든 문자열을 다음 규칙으로 정렬하여 출력하라.
\begin{enumerate}
  \item 길이가 짧은 문자열이 먼저 온다.
  \item 길이가 같다면 사전순으로 빠른 문자열이 먼저 온다.
\end{enumerate}
같은 문자열이 여러 번 주어져도 중복을 제거하지 않는다.

\textbf{구현 및 제출 조건}
\begin{itemize}
  \item C++17 또는 C++20로 완전한 프로그램을 작성한다.
  \item 정렬 규칙이 항상 만족되는 이유를 설명한다.
  \item 시간복잡도와 공간복잡도를 설명한다.
  \item 주요 edge case를 제시한다.
  \item 직접 측정한 풀이 시간을 기록한다.
\end{itemize}

\textbf{제약 조건}
\begin{TextBlock}
1 <= N <= 100000
1 <= 각 문자열의 길이 <= 30
문자열은 영어 소문자로만 이루어진다.
\end{TextBlock}
\end{problembox}

\begin{answerbox}[학습자 제출 답안]
\begin{lstlisting}
#include <iostream>
#include <string>
#include <algorithm>
#include <vector>

using namespace std;

bool cmp(const string& s1, const string& s2) {
    if (s1.size() != s2.size()) {
        return s1.size() < s2.size();
    }
    return s1 < s2;
}

int main()
{
    int N = 0;
    cin >> N;
    vector<string> v = {};
    for (int i = 1; i <= N; i += 1) {
        string s = "str";
        cin >> s;
        v.push_back(s);
    }
    sort(v.begin(), v.end(), cmp);
    for (int j = 0; j < v.size(); j += 1) {
        cout << v[j] << '\n';
    }
}
\end{lstlisting}

\textbf{정렬 설명}: 두 문자열의 길이가 다르면 짧은 문자열을 먼저 두고, 길이가 같으면 사전순으로 빠른 문자열을 먼저 둔다.

\textbf{시간복잡도}: 입력 및 저장은 $O(N)$, 정렬은 $O(N\log N)$이므로 전체는 $O(N\log N)$이다. 문제에서 문자열 길이가 최대 30으로 제한되어 있으므로 문자열 비교 비용은 $N$에 대한 asymptotic 표기에서 상수로 볼 수 있다.

\textbf{공간복잡도}: 최대 길이 30인 문자열을 $N$개 저장하므로 문제의 제한 아래에서 $O(N)$이다.

\textbf{주요 edge case}: $N=1$, 중복 문자열이 존재하는 경우.

\textbf{T\_solve}: 10분 17초.
\end{answerbox}

% =========================================================
% Q&A after Part E
% =========================================================
\section{중간평가 이후 학습 질의응답}

\subsection{Q1. comparator의 argument에서 \code{const}와 \code{\&}를 빼면 어떻게 되는가?}
학습자는 comparator를 다음처럼 작성했을 때 \code{const}와 \code{\&}의 의미를 질문했다.
\begin{lstlisting}
bool cmp(const string& s1, const string& s2)
\end{lstlisting}

비교하면 다음과 같다.
\begin{itemize}
  \item \code{bool cmp(string s1, string s2)}: 각 비교 호출에서 문자열을 복사한다.
  \item \code{bool cmp(string\& s1, string\& s2)}: 복사는 피하지만 원본을 수정할 수 있는 interface가 된다.
  \item \code{bool cmp(const string s1, const string s2)}: 여전히 복사하며, 복사본만 수정하지 못한다.
  \item \code{bool cmp(const string\& s1, const string\& s2)}: 복사하지 않고 읽기만 한다.
\end{itemize}

\code{sort}는 comparator를 여러 번 호출하므로 큰 객체를 값으로 전달하면 불필요한 복사가 반복될 수 있다. 따라서 읽기 전용 비교에는 \code{const T\&}가 자연스럽다.

\subsection{Q2. \code{int j < v.size()} 경고는 왜 생기는가?}
\code{int}는 signed이고 \code{v.size()}의 반환형은 unsigned \code{size\_t} 계열이기 때문이다. 현재 문제처럼 \code{j}가 0에서 시작해 증가하고 $N\le100000$인 경우에는 실제 오류가 발생하지 않지만, compiler는 일반적인 signed/unsigned 변환 위험을 경고한다.

특히 signed 음수를 unsigned와 비교하면 음수가 매우 큰 양수로 변환될 수 있다.

가능한 대안은 다음과 같다.
\begin{lstlisting}
for (size_t j = 0; j < v.size(); ++j) {
    cout << v[j] << '\n';
}
\end{lstlisting}

index가 필요 없다면 다음이 더 단순하다.
\begin{lstlisting}
for (const string& s : v) {
    cout << s << '\n';
}
\end{lstlisting}

\subsection{Q3. 정수를 표현하는 타입에는 무엇이 있는가?}
학습자는 \code{int}, \code{long long}, \code{size\_t}의 관계를 질문했다. 이 셋만 있는 것은 아니며, C++에는 \code{short}, \code{int}, \code{long}, \code{long long}과 각각의 unsigned 형태가 있다.

CP에서는 다음처럼 단순화해 기억할 수 있다.
\begin{itemize}
  \item 일반적인 작은 정수: \code{int}
  \item 계산 결과가 약 $\pm2\times10^9$ 범위를 넘을 수 있음: \code{long long}
  \item container의 크기와 직접 맞춰 쓰는 크기/index: \code{size\_t}
\end{itemize}

다만 값이 양수라는 이유만으로 무조건 \code{size\_t}를 선택하는 것은 아니다. 점수나 가격처럼 개념적으로 정수인 값은 양수 범위라도 \code{int}를 사용할 수 있다.

\subsection{Q4. \code{v.begin()}과 \code{v.end()}는 원소인가, index인가?}
둘 다 아니다. iterator다. \code{v.begin()}은 첫 원소의 위치, \code{v.end()}는 마지막 원소 바로 다음 위치를 나타낸다.

\begin{lstlisting}
auto it = v.begin();
cout << *it;  // iterator가 가리키는 실제 원소
\end{lstlisting}

\code{sort(v.begin(), v.end())}는 $[\code{begin()},\code{end()})$ 범위, 즉 vector 전체를 정렬한다.

\subsection{Q5. \code{auto}는 무엇인가?}
\code{auto}는 초기값을 보고 compiler가 타입을 추론하도록 하는 키워드다.
\begin{lstlisting}
auto x = 10;          // int
auto it = v.begin();  // vector<int>::iterator 계열
\end{lstlisting}

동적 타입을 의미하지 않는다. C++는 여전히 정적 타입 언어이며, compile time에 타입이 결정된다.

또한 reference를 유지하려면 직접 \code{\&}를 써야 한다.
\begin{lstlisting}
for (auto x : v) { }        // copy
for (auto& x : v) { }       // reference
for (const auto& x : v) { } // read-only reference
\end{lstlisting}

\subsection{Q6. \code{int\& r = x;}에서 새로 정의되는 \code{r}이 ``복사되지 않는다''는 말은 무엇인가?}
``\code{r}을 복사하지 않는다''는 뜻이 아니다. \code{x}의 값을 복사해 새로운 \code{int} 객체를 만들지 않는다는 뜻이다.

\begin{lstlisting}
int x = 10;
int& r = x;
\end{lstlisting}

여기서 \code{x}와 \code{r}은 같은 \code{int} 객체를 부르는 두 이름이다. 따라서 \code{r=20;}을 실행하면 \code{x}도 20이 된다.

\subsection{Q7. 함수의 \code{int\& x}, \code{int x}, \code{const int x}는 어떻게 다른가?}
\begin{lstlisting}
void func1(int& x) { ... }
void func2(int x) { ... }
void func3(const int x) { ... }
\end{lstlisting}

자연어로 옮기면 다음과 같다.
\begin{itemize}
  \item \code{int\& x}: 호출자의 원본 \code{int}를 직접 참조하고 수정할 수 있다.
  \item \code{int x}: 값을 복사해 지역 변수로 사용한다. 원본은 바뀌지 않는다.
  \item \code{const int x}: 값을 복사하되 함수 안에서 그 복사본도 수정하지 않는다.
\end{itemize}

큰 container를 읽기만 할 때 자주 쓰는 형태는 \code{const T\&}이다.

\subsection{Q8. 일반 reference 선언과 reference parameter는 엄밀히 다른 의미인가?}
\code{int\& a}라는 타입 의미 자체는 같다. 차이는 bind되는 시점이다.
\begin{lstlisting}
int& a = x;       // declaration time: a binds to x
\end{lstlisting}

반면
\begin{lstlisting}
void func(int& a) {
    // a binds to the argument of each call
}
\end{lstlisting}
은 함수 호출 시 전달된 argument에 bind된다.

% =========================================================
% Part F
% =========================================================
\section{Part F --- 최종 성취도 평가 문제}

\subsection{F-1. Priority List}
\begin{problembox}{최종평가 문제 F-1 --- Priority List}
$N$개의 레코드가 주어진다. 각 레코드는 문자열 \code{name}과 정수 \code{priority}로 구성된다.

다음 우선순위로 정렬한다.
\begin{enumerate}
  \item \code{priority}가 큰 레코드가 먼저 온다.
  \item \code{priority}가 같다면 \code{name}의 길이가 짧은 레코드가 먼저 온다.
  \item 두 조건도 같다면 \code{name}이 사전순으로 빠른 레코드가 먼저 온다.
\end{enumerate}

먼저 정렬된 전체 목록을 출력하고, 이어서 그 목록의 정확한 역순을 출력한다. 동일한 \code{(name, priority)} 레코드가 여러 번 입력되어도 모두 보존한다.

\textbf{구현 및 제출 조건}
\begin{itemize}
  \item C++17 또는 C++20을 사용한다.
  \item 출력 작업은 별도의 함수 \code{print\_records}에서 수행해야 한다.
  \item \code{print\_records}는 출력할 레코드 container 전체를 매개변수로 받아야 한다.
  \item \code{print\_records}는 호출자의 container를 수정해서는 안 된다.
  \item \code{print\_records} 호출 때문에 container 전체가 복사되어서는 안 된다.
  \item 역순 출력을 위해 원본 레코드 전체를 복제한 새로운 $O(N)$ 크기 container를 만들어서는 안 된다.
  \item 중복 레코드를 제거해서는 안 된다.
  \item 완전한 코드, 정렬 이유, 함수 매개변수 선택 이유, 시간복잡도, 공간복잡도, edge case, 직접 측정한 풀이 시간을 제출한다.
\end{itemize}

\textbf{제약 조건}
\begin{TextBlock}
1 <= N <= 100000
1 <= name의 길이 <= 30
name은 영어 소문자로만 이루어진다.
-1000000000 <= priority <= 1000000000
\end{TextBlock}

\textbf{예제 입력}
\begin{TextBlock}
7
pear 3
a 5
kiwi 3
fig 5
plum 3
fig 5
banana 3
\end{TextBlock}

\textbf{예제 출력}
\begin{TextBlock}
FORWARD
a 5
fig 5
fig 5
kiwi 3
pear 3
plum 3
banana 3
REVERSED
banana 3
plum 3
pear 3
kiwi 3
fig 5
fig 5
a 5
\end{TextBlock}
\end{problembox}

\begin{answerbox}[학습자 제출 답안]
\begin{lstlisting}
#include <iostream>
#include <string>
#include <algorithm>
#include <vector>

using namespace std;

bool print_records(const pair<string, int>& p1,
                   const pair<string, int>& p2) {
    if (p1.second != p2.second) {
        return p1.second > p2.second;
    }
    else if (p1.second == p2.second &&
             p1.first.size() != p2.first.size()) {
        return p1.first.size() < p2.first.size();
    }
    else if (p1.second == p2.second &&
             p1.first.size() == p2.first.size()) {
        return p1.first < p2.first;
    }
}

int main()
{
    int N = 0;
    cin >> N;
    vector<pair<string, int>> v = {};
    for (int i = 1; i <= N; i += 1) {
        string n = "name";
        int p = 0;
        cin >> n >> p;
        pair<string, int> P = { n, p };
        v.push_back(P);
    }
    cout << "FORWARD" << '\n';
    sort(v.begin(), v.end(), print_records);
    for (int i = 0; i < v.size(); i += 1) {
        string name = v[i].first;
        int pri = v[i].second;
        cout << name << " " << pri << '\n';
    }
    cout << "REVERSED" << '\n';
    reverse(v.begin(), v.end());
    for (int i = 0; i < v.size(); i += 1) {
        string name = v[i].first;
        int pri = v[i].second;
        cout << name << " " << pri << '\n';
    }
}
\end{lstlisting}

\textbf{정렬 설명}: priority가 다르면 priority만 비교하고, 같으면 이름 길이를 비교하며, 그것도 같으면 이름을 사전순으로 비교한다.

\textbf{매개변수 설명}: 레코드는 \code{pair<string,int>}이고, 레코드를 복사하지 않고 수정하지 않기 위해 comparator의 두 매개변수를 \code{const \&}로 받았다.

\textbf{시간복잡도}: 선형 입력/출력 및 \code{reverse}는 $O(N)$, \code{sort}는 $O(N\log N)$이므로 전체 $O(N\log N)$.

\textbf{공간복잡도}: 입력 record 저장에 $O(N)$.

\textbf{주요 edge case}: $N=1$, 중복 record와 여러 priority가 섞인 경우.

\textbf{T\_solve}: 24분 20초.
\end{answerbox}

\subsection{개념 보강: comparator와 container 처리 함수의 역할 구분}
두 함수는 모두 \code{const \&}를 사용할 수 있지만 역할이 다르다.

\begin{lstlisting}
bool cmp(const pair<string, int>& a,
         const pair<string, int>& b)
{
    // compare two records
}

void print_records(
    const vector<pair<string, int>>& records)
{
    // read and print the whole container
}
\end{lstlisting}

첫 번째 함수는 \textbf{레코드 두 개}의 순서를 판단한다. 두 번째 함수는 \textbf{레코드 전체 container}를 전달받아 처리한다. 문제의 함수 interface 조건을 읽을 때 ``함수 이름''뿐 아니라 ``함수가 받아야 하는 대상과 책임''을 함께 확인해야 한다.

\subsection{F-2. Item Ordering}
\begin{problembox}{최종평가 재평가 문제 F-2 --- Item Ordering}
$N$개의 상품 record가 주어진다. 각 record는 문자열 \code{name}과 정수 \code{cost}로 구성된다.

모든 record를 다음 순서로 정렬한다.
\begin{enumerate}
  \item \code{cost}가 작은 record가 먼저 온다.
  \item \code{cost}가 같다면 \code{name} 길이가 긴 record가 먼저 온다.
  \item 두 조건도 같다면 \code{name}이 사전순으로 빠른 record가 먼저 온다.
\end{enumerate}

정렬된 전체 목록을 출력한 뒤, 그 목록의 정확한 역순을 다시 출력한다. 동일한 \code{(name,cost)} record도 모두 유지한다.

\textbf{구현 및 제출 조건}
\begin{itemize}
  \item C++17 또는 C++20을 사용한다.
  \item record는 \code{vector<pair<string,int>>}에 저장한다.
  \item 출력은 반드시 별도의 함수 \code{write\_list}에서 수행한다.
  \item \code{write\_list}는 전체 \code{vector<pair<string,int>>}를 하나의 매개변수로 받아야 한다.
  \item 호출자의 container를 수정해서는 안 된다.
  \item 호출할 때 container 전체가 복사되어서는 안 된다.
  \item 역순 출력을 위해 원본과 같은 크기의 새로운 $O(N)$ record container를 만들면 안 된다.
  \item 중복 record를 삭제하면 안 된다.
  \item 완전한 코드, 정렬 이유, 함수 매개변수 선택 이유, 시간복잡도, 공간복잡도, edge case, 직접 측정한 풀이 시간을 제출한다.
\end{itemize}

\textbf{제약 조건}
\begin{TextBlock}
1 <= N <= 100000
1 <= name의 길이 <= 30
name은 영어 소문자로만 구성된다.
-1000000000 <= cost <= 1000000000
\end{TextBlock}

\textbf{예제 입력}
\begin{TextBlock}
8
kiwi 5
a 2
pear 5
fig 2
plum 5
banana 5
fig 2
zz 5
\end{TextBlock}

\textbf{예제 출력}
\begin{TextBlock}
NORMAL
fig 2
fig 2
a 2
banana 5
kiwi 5
pear 5
plum 5
zz 5
BACKWARD
zz 5
plum 5
pear 5
kiwi 5
banana 5
a 2
fig 2
fig 2
\end{TextBlock}
\end{problembox}

\begin{answerbox}[학습자 제출 답안]
\begin{lstlisting}
#include <iostream>
#include <string>
#include <algorithm>
#include <vector>

using namespace std;

bool cmp(const pair<string, int>& r1,
         const pair<string, int>& r2)
{
    string name1 = r1.first;
    int cost1 = r1.second;

    string name2 = r2.first;
    int cost2 = r2.second;

    if (cost1 != cost2) {
        return cost1 < cost2;
    }
    else if (cost1 == cost2 &&
             name1.size() != name2.size()) {
        return name1.size() > name2.size();
    }
    else if (cost1 == cost2 &&
             name1.size() == name2.size()) {
        return name1 < name2;
    }
    else {

    }
}

void write_list(const vector<pair<string, int>>& v) {
    cout << "NORMAL" << '\n';
    for (int i = 0; i < v.size(); i += 1) {
        string name = v[i].first;
        int cost = v[i].second;
        cout << name << " " << cost << '\n';
    }
    cout << "BACKWARD" << '\n';
    for (int i = 1; i <= v.size(); i += 1) {
        string name = v[v.size() - i].first;
        int cost = v[v.size() - i].second;
        cout << name << " " << cost << '\n';
    }
}

int main()
{
    int N = 0;
    vector<pair<string, int>> v = {};
    cin >> N;
    for (int i = 1; i <= N; i += 1) {
        string name = "n";
        int cost = 0;
        cin >> name >> cost;
        pair<string, int> p = {name, cost};
        v.push_back(p);
    }
    sort(v.begin(), v.end(), cmp);
    write_list(v);
}
\end{lstlisting}

\textbf{정렬 설명}: cost가 다르면 cost 오름차순으로 정렬하고, 같으면 이름 길이 내림차순, 다시 같으면 사전순으로 정렬한다.

\textbf{매개변수 설명}: 여러 상품을 담은 \code{vector<pair<string,int>>}를 복사하거나 수정할 필요가 없으므로 \code{const \&}로 받는다.

\textbf{시간복잡도}: 입력과 출력은 $O(N)$, 정렬은 $O(N\log N)$이므로 전체 $O(N\log N)$.

\textbf{공간복잡도}: 입력 record를 저장하는 vector가 $O(N)$.

\textbf{주요 edge case}: $N=1$, 같은 이름/다른 cost, 중복 record가 있는 경우.

\textbf{T\_solve}: 20분 51초.
\end{answerbox}

\subsection{코드 품질 보강: 논리적으로 exhaustive여도 return을 명시하기}
위 comparator에서는 세 조건이 논리적으로 가능한 모든 경우를 덮지만, 함수 끝에 syntactic return이 없는 형태라 compiler가 ``control reaches end of non-void function'' 경고를 낼 수 있다.

다음처럼 마지막 경우를 직접 return으로 두면 더 명확하다.
\begin{lstlisting}
bool cmp(const pair<string, int>& r1,
         const pair<string, int>& r2)
{
    if (r1.second != r2.second) {
        return r1.second < r2.second;
    }
    if (r1.first.size() != r2.first.size()) {
        return r1.first.size() > r2.first.size();
    }
    return r1.first < r2.first;
}
\end{lstlisting}

또한 comparator 내부에서
\begin{lstlisting}
string name1 = r1.first;
string name2 = r2.first;
\end{lstlisting}
처럼 문자열을 다시 복사할 필요는 없다. 이미 매개변수를 \code{const \&}로 받았으므로 \code{r1.first}, \code{r2.first}를 직접 읽으면 된다.

% =========================================================
% Part G
% =========================================================
\section{Part G --- Adaptive Extra Problem}

\subsection{G-1. Record Highs and Longest Nondecreasing Segment}
\begin{problembox}{Adaptive Extra Problem G-1}
정수 수열 \code{v}가 주어진다. 다음 두 값을 구하라.
\begin{enumerate}
  \item 첫 번째 원소를 제외하고, 지금까지 등장한 모든 이전 원소보다 \textbf{엄격하게 큰} 원소의 개수
  \item 연속 부분구간 중 값이 감소하지 않는 가장 긴 구간의 길이
\end{enumerate}
두 번째 조건은
\[
\code{v}[i] \le \code{v}[i+1] \le \cdots \le \code{v}[j]
\]
를 만족하는 가장 긴 연속 구간의 길이를 뜻한다.

\textbf{구현 및 제출 조건}
\begin{itemize}
  \item C++17 또는 C++20을 사용한다.
  \item 입력 수열은 \code{vector<int>}에 저장한다.
  \item 분석은 반드시 별도의 함수 \code{analyze}에서 수행한다.
  \item \code{analyze}는 전체 vector를 매개변수로 받아야 한다.
  \item 호출자의 vector를 수정해서는 안 된다.
  \item 호출로 vector 전체가 복사되어서는 안 된다.
  \item 반환형은 반드시 \code{pair<int,int>}여야 한다.
  \item 반환값의 첫 번째 성분은 record high 개수, 두 번째 성분은 가장 긴 nondecreasing 연속 구간 길이다.
  \item 전체 알고리즘의 시간복잡도는 $O(N)$이어야 한다.
  \item 원본 vector 외에 $O(N)$ 크기의 추가 container를 만들면 안 된다.
  \item 완전한 코드, 두 계산의 correctness 설명, 매개변수 선택 이유, 시간복잡도, 공간복잡도와 edge case, 직접 측정한 풀이 시간을 제출한다.
\end{itemize}

\textbf{제약 조건}
\begin{TextBlock}
1 <= N <= 100000
-1000000000 <= v[i] <= 1000000000
\end{TextBlock}

\textbf{예제 입력}
\begin{TextBlock}
8
3 3 5 4 4 6 7 2
\end{TextBlock}

\textbf{예제 출력}
\begin{TextBlock}
3 4
\end{TextBlock}
\end{problembox}

\begin{answerbox}[학습자 제출 답안]
\begin{lstlisting}
#include <iostream>
#include <string>
#include <algorithm>
#include <vector>

using namespace std;

pair<int, int> analyze(const vector<int>& v) {
    pair<int, int> result = {};
    int record_count = -1;
    int longest_length = 0;
    int temp_length = 0;
    int max = -1e9-1;
    int previous = -1e9 - 1;

    for (int i = 0; i < v.size(); i += 1) {
        int current = v[i];
        if (v[i] > max) {
            record_count += 1;
            max = v[i];
        }
        if (v[i] >= previous) {
            temp_length += 1;

            if (temp_length > longest_length) {
                longest_length = temp_length;
            }
        }
        else {
            temp_length = 1;
        }
        previous = v[i];
    }

    result = { record_count, longest_length };

    return result;
}

int main()
{
    int N = 0;
    vector<int> v = {};
    cin >> N;
    for (int i = 1; i <= N; i += 1) {
        int n = 0;
        cin >> n;
        v.push_back(n);
    }
    pair<int, int> ans = analyze(v);
    cout << ans.first << " " << ans.second << '\n';
}
\end{lstlisting}

\textbf{record high 설명}: 현재 원소가 이전 최대값보다 클 때만 최대값을 갱신하고 count를 증가시킨다. 첫 원소를 record high 개수에서 제외하기 위해 초기 count를 \code{-1}로 두었으며, 초기 최대값은 문제의 가능한 최솟값보다 작게 둔다.

\textbf{longest nondecreasing 설명}: 현재 원소가 바로 전 원소 이상이면 현재 구간 길이를 1 증가시키고, 작아지면 현재 원소 하나에서 새 구간을 시작하므로 길이를 1로 만든다. 가장 큰 현재 구간 길이를 별도 변수에 저장한다.

\textbf{매개변수 설명}: 입력 vector를 복사하거나 수정할 필요가 없으므로 \code{const \&}로 받는다.

\textbf{시간복잡도}: vector를 한 번 순회하며 각 iteration에서 상수 시간 연산만 하므로 $O(N)$.

\textbf{공간복잡도}: 입력 vector 저장에 $O(N)$, 입력 vector를 제외한 추가 공간은 scalar와 pair뿐이므로 $O(1)$.

\textbf{주요 edge case}: $N=1$, \code{v=\{12789\}}일 때 출력은 \code{0 1}.

\textbf{T\_solve}: 24분 21초.
\end{answerbox}

\subsection{개념 보강: 문제 제한보다 작은 sentinel과 상태 초기화}
학습자 코드는 초기값을 다음처럼 두었다.
\begin{lstlisting}
int max = -1e9 - 1;
int previous = -1e9 - 1;
\end{lstlisting}

문제의 최솟값이 $-10^9$이므로 입력보다 작은 sentinel을 둘 수 있다. 다만 \code{1e9}는 floating-point literal이므로, 정수 범위를 명확하게 표현하려면 다음처럼 정수 literal을 직접 쓰는 방식이 더 읽기 쉽다.
\begin{lstlisting}
int maximum = -1000000001;
int previous = -1000000001;
\end{lstlisting}

또는 첫 원소로 상태를 초기화한 뒤 index 1부터 순회하는 구조도 가능하다. 어느 방식을 사용하든 핵심은 \textbf{초기 상태의 의미와 첫 iteration의 동작을 명확히 설명할 수 있는가}이다.

% =========================================================
% Part H
% =========================================================
\section{Part H --- 숙달 점검과 복습}

\subsection{H1. 개념 자기점검}
다음 질문에 자료 없이 답할 수 있는지 확인하라.
\begin{enumerate}
  \item \code{vector}와 \code{array}를 선택하는 가장 기본적인 기준은 무엇인가?
  \item \code{v.begin()}과 \code{v.end()}는 원소, index, iterator 중 무엇인가?
  \item 왜 \code{*v.end()}는 유효한 원소 접근이 아닌가?
  \item STL 범위 $[\text{first},\text{last})$에서 마지막 iterator가 제외되는 이유를 설명할 수 있는가?
  \item \code{auto}는 동적 타입인가? 아니라면 정확히 무엇인가?
  \item \code{for (auto x : v)}, \code{for (auto\& x : v)}, \code{for (const auto\& x : v)}의 차이는 무엇인가?
  \item \code{int r = x;}와 \code{int\& r = x;}는 객체 관점에서 어떻게 다른가?
  \item \code{T x}, \code{T\& x}, \code{const T\& x} 함수 매개변수의 차이를 설명할 수 있는가?
  \item comparator의 \code{cmp(a,b)}가 참이라는 것은 무엇을 뜻하는가?
  \item comparator에서 \code{<=}를 쓰면 왜 문제가 될 수 있는가?
  \item \code{sort}의 시간복잡도는 무엇인가?
  \item \code{int i < v.size()}에서 경고가 생길 수 있는 이유는 무엇인가?
  \item unsigned index로 역순 순회할 때 어떤 underflow 위험이 있는가?
  \item \code{pair<string,int>}를 어떤 종류의 문제에서 쓰면 자연스러운가?
  \item container 전체를 읽기만 하는 함수에서 왜 \code{const vector<T>\&}가 자주 사용되는가?
\end{enumerate}

\subsection{H2. 구현 자기점검}
\begin{checklistbox}[S0-B 코드 작성 후 체크리스트]
\begin{itemize}
  \item[$\square$] 입력 크기가 고정인지 가변인지 확인했는가?
  \item[$\square$] index가 정말 필요한가, range-for로 충분한가?
  \item[$\square$] \code{begin()}과 \code{end()}의 범위를 정확히 사용했는가?
  \item[$\square$] \code{end()}를 dereference하지 않았는가?
  \item[$\square$] 기본 정렬과 custom comparator 중 적절한 방식을 선택했는가?
  \item[$\square$] comparator가 strict ordering을 만족하는가?
  \item[$\square$] comparator에 모든 return 경로가 명시되어 있는가?
  \item[$\square$] 큰 문자열/container를 불필요하게 값으로 복사하지 않는가?
  \item[$\square$] 읽기 전용 큰 객체에 \code{const \&}를 고려했는가?
  \item[$\square$] 문제에서 요구한 함수 이름뿐 아니라 함수의 역할과 parameter 대상까지 확인했는가?
  \item[$\square$] signed/unsigned 비교 경고가 있는가?
  \item[$\square$] unsigned 역순 순회가 underflow하지 않는가?
  \item[$\square$] 중복 record를 보존해야 하는 문제에서 실수로 제거하지 않았는가?
  \item[$\square$] 전체 저장 공간과 auxiliary space를 구분해서 설명할 수 있는가?
\end{itemize}
\end{checklistbox}

\subsection{H3. S0-B에서 형성해야 할 핵심 사고}
S0-B의 핵심은 STL API 이름을 외우는 것이 아니다. 다음 질문을 코드 작성 전에 스스로 던지는 습관이 중요하다.
\begin{enumerate}
  \item 데이터를 어떤 단위로 저장해야 하는가?
  \item container의 크기와 변경 가능성은 어떤가?
  \item 순회에서 index가 필요한가?
  \item 정렬 순서는 기본 순서인가, 복합 우선순위인가?
  \item 함수에 값을 복사할 필요가 있는가?
  \item 원본 수정이 필요한가, 읽기만 하면 되는가?
  \item 문제의 interface 조건이 구현 구조에 정확히 반영되었는가?
\end{enumerate}

\subsection{H4. 다음 Unit과의 연결}
다음 Learning Unit인 S0-C에서는 이번 Unit에서 이미 사용한 복잡도와 정수형 판단을 더 체계적으로 다룬다.
\begin{itemize}
  \item $O(N)$, $O(N\log N)$, $O(N^2)$의 성장 차이
  \item 입력 크기로부터 실행 가능성을 추정하는 방법
  \item 전체 공간복잡도와 auxiliary space의 구분
  \item \code{int}, \code{long long}의 범위
  \item overflow가 연산 중 어느 시점에 발생하는가
  \item 자료형 선택을 문제의 최대 가능한 값에서 역산하는 방법
\end{itemize}

